\documentclass[10pt,conference]{IEEEtran}
\usepackage{booktabs}
\usepackage{url}
\usepackage{graphicx}

\title{Benevente Quant AI: An Auditable Value--Quality Research and Shadow-Portfolio Framework for Brazilian Equities}
\author{Author information withheld for double-blind review}

\begin{document}
\maketitle

\begin{abstract}
This paper presents Benevente Quant AI, a reproducible research framework for Brazilian-equity and CDI allocation. The framework separates point-in-time data controls, deterministic value--quality eligibility, optional structured large-language-model review, and constrained portfolio optimization. An archived evaluation over pre-specified 5-, 10-, and 15-year windows applies modeled frictions and compares the strategy with CDI, Ibovespa, and constrained mean--variance optimization. The strategy did not exceed CDI over five years, exceeded CDI in the reported 10- and 15-year windows, and was slightly below classical constrained optimization over fifteen years; it therefore provides conditional historical evidence, not a general outperformance claim. The operational component is a R\$100,000 prospective shadow-portfolio protocol. It uses CVM ITR and DFP filings, attributable market-data snapshots, archived price history, liquidity and cost gates, lot-rounded review-only orders, and broker-note reconciliation. No broker API or autonomous investment action is enabled.
\end{abstract}

\begin{IEEEkeywords}
portfolio optimization, Brazilian equities, CVM, reproducibility, value investing, audit trail, large language models
\end{IEEEkeywords}

\section{Introduction}
Long-horizon allocation requires more than an optimizer: returns, accounting reports, liquidity, costs, constraints, and the information available on a decision date must agree. Benevente Quant AI is a research framework designed to make that process inspectable. Its commercial B2B presentation is called Benevente Wealth System; this paper addresses the academic framework.

The system is not an autonomous trading engine and does not claim to beat CDI, Ibovespa, or any other benchmark. Its research question is narrower: whether a transparent, constrained process can be evaluated reproducibly after modeled frictions. Individualized securities advice and portfolio management for third parties require the applicable legal, contractual, and professional structure; an algorithm does not remove those responsibilities \cite{cvm19,cvmrobo}.

The framework follows Markowitz's risk--return setting while treating estimation error, concentration and turnover as first-class constraints. Black and Litterman's equilibrium-plus-views formulation motivates a strict separation between qualitative views and mechanical allocation \cite{blacklitterman}; the present implementation does not claim to be a Black--Litterman allocator. Evidence that optimized portfolios often fail to dominate simple diversification out of sample motivates explicit caps and a frozen benchmark protocol \cite{demiguel}. Transaction frictions are evaluated as part of the strategy rather than appended as an afterthought, consistent with evidence that costs and turnover can materially change anomaly profitability \cite{novymarx}.

\section{Method}
At every rebalance, estimators only use returns available through $t-1$. The portfolio weights $w$ solve a long-only constrained mean--variance problem:
\begin{equation}
 \max_w \quad \hat{\mu}^{\top}w - \frac{\gamma}{2}w^{\top}\hat{\Sigma}w,
 \quad \sum_i w_i=1,
\end{equation}
subject to non-negative weights, an issuer cap, a policy-specific equity cap, and a residual CDI sleeve. The optimization is deterministic once inputs and policy are fixed.

Before optimization, the value--quality selector rejects securities with insufficient capitalization or liquidity. For non-financial issuers it requires positive free-cash-flow yield, minimum ROIC, verified debt-to-EBITDA, and verified interest coverage. For financial issuers it uses metrics suited to their statements, including ROE and price-to-book. Missing or non-comparable fields are rejection conditions. An optional LLM can create a structured skeptical thesis for assets already eligible, but cannot select an asset, define weights, override a hard filter, or submit an order.

Historical experiments apply 10 basis points of transaction cost plus 5 basis points of slippage to executed turnover. The live proposal separately estimates Clear/B3 assumptions and requires reconciliation to the actual broker note. These are assumptions, not observed universal trading costs.

\section{Data and Point-in-Time Controls}
The archived backtest uses adjusted B3 prices and Ibovespa data obtained through Yahoo Finance, with CDI, Selic, and IPCA from Banco Central do Brasil SGS. Inputs are stored with each run. The current universe is fixed and based on surviving tickers, so survivorship bias remains a material limitation.

For a current proposal, Benevente downloads CVM quarterly ITR and annual DFP archives \cite{cvmitr,cvmdfp}. For income-statement and cash-flow items, it computes
\begin{equation}
 X_{TTM}=X_{DFP,y-1}+X_{ITR,q,y}-X_{ITR,q,y-1}.
\end{equation}
The receipt date $DT\_RECEB$ is the availability date and $DT\_REFER$ is the accounting observation date. A filing received after the decision date is rejected. Market capitalization, average daily traded value, closing price, and lot size come from a dated B3, broker, or licensed-vendor snapshot. The optimizer uses a separately archived price-history CSV, including a CDI sleeve; future observations, null prices, or an insufficient history block the run.

The investor horizon is selectable at 1, 2, 5, 10, or 15 years. It changes the return-estimation sample to 252, 504, 756, 1260, or 1260 trading days, respectively. This makes the horizon an explicit estimation choice, while risk profile controls the equity cap, issuer cap, drawdown tolerance, and review interval.

\section{Archived Empirical Evaluation}
The 5-, 10-, and 15-year windows were pre-specified to end on 30 June 2026 and start on 1 July 2021, 2016, and 2011. Each uses a separate one-year lookback. Independent validation checks input continuity, duplicate dates, null values, and wealth-curve recomposition.

\begin{table*}[t]
\caption{Archived horizon results after modeled frictions.}
\centering
\begin{tabular}{llrrrr}
\toprule
Horizon & Strategy & Cumulative return & CAGR & Annual volatility & Maximum drawdown \\
\midrule
5 years & Benevente Quant AI & 69.87\% & 11.38\% & 8.69\% & -11.00\% \\
        & CDI & 76.47\% & 12.25\% & 0.67\% & 0.00\% \\
        & Ibovespa & 38.09\% & 6.78\% & 16.78\% & -20.43\% \\
10 years & Benevente Quant AI & 226.23\% & 12.78\% & 15.04\% & -29.59\% \\
         & CDI & 141.66\% & 9.39\% & 1.13\% & 0.00\% \\
         & Ibovespa & 220.60\% & 12.58\% & 24.79\% & -45.68\% \\
15 years & Benevente Quant AI & 408.29\% & 11.65\% & 12.23\% & -18.93\% \\
         & CDI & 298.23\% & 9.82\% & 1.00\% & 0.00\% \\
         & Ibovespa & 174.10\% & 7.08\% & 23.39\% & -40.77\% \\
\bottomrule
\end{tabular}
\label{tab:horizons}
\end{table*}

The results are mixed. Benevente Quant AI did not beat CDI over five years. It exceeded CDI in the reported ten- and fifteen-year windows; in the ten-year window it had the largest cumulative return among the displayed benchmarks. Over fifteen years, classical constrained MVO returned 413.78\%, modestly above Benevente's 408.29\%. These observations do not establish persistent alpha, causal benefit from the LLM, or a performance expectation.

\section{Prospective R\$100,000 Shadow Portfolio}
The prospective protocol starts with a R\$100,000 virtual portfolio, not automatic investment. A policy records a risk profile (conservative, moderate, growth, aggressive, or custom), horizon, equity and issuer limits, drawdown tolerance, maximum rebalance cost, review interval, and excluded tickers. It must be explicitly acknowledged before a proposal is issued.

The proposal gate requires: (i) the most recent CVM filing that was public at the decision date; (ii) TTM fundamentals; (iii) market and price-history files that are dated and attributable; (iv) a fail-closed fundamental screen; and (v) concentration, freshness, liquidity, and cost checks. The package preserves the policy, snapshots, fundamentals, screen, weights, data-source metadata, and order summary. Initial buy instructions are rounded down to the declared lot, include a limit reference price and estimated cost, and are rejected if an order would exceed 5\% of average daily traded value. The system does not send them anywhere.

If a responsible person elects to act, they independently approve each line and enter it manually at the broker. An execution export with broker-note identifier, quantity, price, and fees is subsequently reconciled against the proposed order. Prospective NAV reporting separately records portfolio value, CDI, Ibovespa, realized costs, turnover, and drawdown. Twelve, twenty-four, and sixty-month prospective observations should be appended without modifying the frozen historical experiment.

\section{Limitations and Conclusion}
The value--quality extension has no historical return claim yet because a historical point-in-time fundamental panel has not been constructed across the full backtest period. The fixed universe, secondary historical price source, modeled costs, taxes, bid--ask spreads, corporate actions, delistings, and realized broker conditions all limit inference. The shadow portfolio is a protocol for evidence collection, not investment validation.

Benevente Quant AI nevertheless provides an auditable research baseline: data availability is explicit, filters fail closed, constraints are deterministic, performance artifacts can be reproduced, and the current proposal workflow archives data and requires human review. The principal next research task is a historical CVM fundamental panel with dated constituents and filings, evaluated walk-forward without altering rules after results are observed.

\begin{thebibliography}{00}
\bibitem{markowitz} H. Markowitz, ``Portfolio Selection,'' \emph{The Journal of Finance}, vol. 7, no. 1, pp. 77--91, 1952.
\bibitem{blacklitterman} F. Black and R. Litterman, ``Global Portfolio Optimization,'' \emph{Financial Analysts Journal}, vol. 48, no. 5, pp. 28--43, 1992, doi: 10.2469/faj.v48.n5.28.
\bibitem{demiguel} V. DeMiguel, L. Garlappi, and R. Uppal, ``Optimal versus Naive Diversification: How Inefficient is the 1/N Portfolio Strategy?'' \emph{The Review of Financial Studies}, vol. 22, no. 5, pp. 1915--1953, 2009, doi: 10.1093/rfs/hhm075.
\bibitem{novymarx} R. Novy-Marx and M. Velikov, ``A Taxonomy of Anomalies and Their Trading Costs,'' \emph{The Review of Financial Studies}, vol. 29, no. 1, pp. 104--147, 2016, doi: 10.1093/rfs/hhv063.
\bibitem{cvmitr} Comiss\~ao de Valores Mobili\'arios, ``Formul\'ario de Informa\c{c}\~oes Trimestrais (ITR),'' CVM Open Data. [Online]. Available: \url{https://dados.cvm.gov.br/dados/CIA_ABERTA/DOC/ITR/DADOS/}
\bibitem{cvmdfp} Comiss\~ao de Valores Mobili\'arios, ``Demonstra\c{c}\~oes Financeiras Padronizadas (DFP),'' CVM Open Data. [Online]. Available: \url{https://dados.cvm.gov.br/dataset/cia_aberta-doc-dfp}
\bibitem{cvmrobo} Comiss\~ao de Valores Mobili\'arios, ``Rob\^os de investimentos,'' Portal do Investidor. [Online]. Available: \url{https://www.gov.br/investidor/pt-br/investir/como-investir/profissionais-do-mercado/robos-de-investimentos}
\bibitem{cvm19} Comiss\~ao de Valores Mobili\'arios, ``Resolu\c{c}\~ao CVM n. 19,'' 2021. [Online]. Available: \url{https://conteudo.cvm.gov.br/legislacao/resolucoes/resol019.html}
\end{thebibliography}
\end{document}
